\section{Rational descent with closed supports}
\label{sec:descent}

We compare ordinary and normalized Euler lengths through rational
$K$-theory with support. In the regular-base applications, the strategy
is to compare the two evaluations on classes pulled back from the base,
then show that sums of conjugate classes come from rational classes over
that base. The latter step is the descent theorem below; the final
subsection combines it with the normalized-length input of
Section~\ref{sec:lengths}. We retain connective associative coefficients
because Section~\ref{sec:asymptotic} applies the same descent theorem to
the quotient by negligible-rank sequences.

All derived algebras in this section are indexed homologically. Thus
connective means that their homotopy groups vanish in negative degrees.

Let $A$ be a Noetherian commutative $\Z_{(p)}$-algebra of finite Krull
dimension, let $J\subset A$ be an ideal containing $p$, and let
$H$ be a connective associative $A$-algebra. The image of $A$ acts
centrally. For a commutative $A$-algebra $B$, write
\[
 \Perf_J(H\otimes_A^{\mathbf L}B)
\]
for the category of perfect modules whose central localizations vanish
off $V(JB)$. Its nonconnective $K$-theory spectrum and its zeroth
homotopy group are denoted by
$\K^J(H\otimes_A^{\mathbf L}B)$ and
$K_0^J(H\otimes_A^{\mathbf L}B)$, respectively. Rationalization of a
spectrum is indicated by a subscript $\Q$.
For a finite group action, $(-)^{hG}$ denotes homotopy fixed points
(derived invariants), whereas $(-)^G$ denotes ordinary invariants.

\begin{theorem}[Supported rational descent]
\label{ds:descent}
Let $A$ be a commutative Noetherian $\Z_{(p)}$-algebra of finite Krull
dimension, let $J\subset A$ be an ideal containing $p$, and let $H$ be a connective
associative $A$-algebra with central $A$-action.
Let $T$ be a finite commutative $A$-algebra with an action of a finite
group $G$ over $A$. Suppose that $\Spec T\to\Spec A$ is surjective and
that $G$ acts transitively on the points of every geometric fiber.
Then pullback induces an equivalence
\begin{equation}
 \K^J(H)_\Q
 \simeq
 \bigl(\K^J(H\otimes_A^{\mathbf L}T)_\Q\bigr)^{hG}.
 \label{ds:spectral-descent}
\end{equation}
In particular, the map
\begin{equation}
 K_0^J(H)\otimes\Q
 \longrightarrow
 \bigl(K_0^J(H\otimes_A^{\mathbf L}T)\otimes\Q\bigr)^G
 \label{ds:kzero-descent}
\end{equation}
is surjective.
\end{theorem}

Neither flatness nor tame ramification of $A\to T$ is required.
The coefficient algebra $H$ need not be commutative, finite over $A$,
or bounded in positive homological degrees. Its connectivity is
essential to the argument.

\subsection{The truncating invariant and its supports}

We use the following precise form of the results of Land--Tamme.
A localizing invariant takes exact sequences of stable categories to
fiber sequences. It is called \emph{truncating} if its value on a
connective associative algebra is unchanged by passage to $\pi_0$.
Write $\operatorname{fib}$ for homotopy fiber and
$\operatorname{HN}(-/\Q)$ for negative cyclic homology over $\Q$.
Rational infinitesimal $K$-theory
\[
 I(U)=\operatorname{fib}\bigl(
   \K(U)_\Q\longrightarrow
   \operatorname{HN}(U\otimes\Q/\Q)\bigr)
\]
is truncating by Goodwillie's relative comparison theorem
\cite[Main Theorem]{Goodwillie}, with the passage from simplicial rings
to connective associative ring spectra explained in
\cite[proof of Corollary~3.9 and Remark~3.10]{LT}. Truncating localizing
invariants satisfy cdh descent, also after tensoring with a fixed
connective associative algebra over the base
\cite[Theorem~A.2 and Appendix~B, especially Theorem~B.3]{LT}.
The cdh topology is generated by Nisnevich covers and abstract blowup
squares: a proper map and a closed center, with the proper map an
isomorphism off that center.

For a derived $\Z_{(p)}$-algebra $U$, consider
\begin{equation}
 E(U)=\operatorname{fib}\bigl(
      \K(U)_\Q\longrightarrow\K(U[1/p])_\Q\bigr).
 \label{ds:relative-invariant}
\end{equation}
The negative cyclic homology terms in the two character maps are
naturally identical, since
$U\otimes\Q\simeq U[1/p]\otimes\Q$. Taking fibers of the character
square therefore gives
\begin{equation}
 E(U)\simeq\operatorname{fib}\bigl(I(U)\longrightarrow I(U[1/p])\bigr).
 \label{ds:character-cancellation}
\end{equation}
It follows that $E$ is a truncating localizing invariant of
$\Z_{(p)}$-linear categories. Notice that this argument does not
identify $\K(U)_\Q$ with $\K(U\otimes\Q)$.

Tensoring with $H$ preserves the localizing property. It also
preserves truncation: if $B$ is connective, then
$H\otimes_A^{\mathbf L}B$ is connective and
\[
 \pi_0(H\otimes_A^{\mathbf L}B)
   =\pi_0(H)\otimes_A\pi_0(B).
\]
Thus replacing $B$ by $\pi_0(B)$ does not change the value of $E$.
For a finite-type $A$-scheme $Y$, let $H_Y$ be the derived pullback
of $H$ and put $Z_Y=Y\times_A V(J)$. Define
\begin{equation}
 \cL(Y)=\operatorname{fib}\bigl(
    E(Y;H)\longrightarrow E(Y\setminus Z_Y;H)\bigr).
 \label{ds:supported-functor}
\end{equation}
Here the coefficient notation means evaluation on the category of
perfect $H_Y$-modules. The coefficient theorem just cited gives cdh
descent for $E(-;H)$. Restriction to the open complement of $Z_Y$
preserves the squares defining this topology, and limits commute
with limits. Hence $\cL$ also satisfies cdh descent.

Because $p\in J$, the support $Z_Y$ disappears after inverting $p$.
Expanding the two fibers in \eqref{ds:supported-functor} and using
perfect-complex localization consequently identifies
\begin{equation}
 \cL(Y)\simeq\K\bigl(\Perf_{Z_Y}(H_Y)\bigr)_\Q.
 \label{ds:support-identification}
\end{equation}
In particular, this is the desired rational supported $K$-theory,
not a cdh replacement of it. The functor is nilinvariant as well.

We recall why the support condition is compatible with the
coefficient categories used here. If $P$ is perfect and $z\in A$,
compactness identifies localization of its endomorphism spectrum
with the endomorphism spectrum of $P[1/z]$. Therefore
$P[1/z]=0$ if and only if some power $z^v$ kills
$\operatorname{id}_P$ in the homotopy category. Applying this to
finitely many generators of $J$ gives the central support
criterion. The corresponding Koszul tensors generate the compact
kernel of restriction to the complement. Perfect localization,
including idempotent completion, thus computes the fiber in
\eqref{ds:support-identification}.

The geometric localization input is Thomason--Trobaugh's perfect
localization theorem \cite[Propositions~5.2.2--5.2.4 and
Theorem~7.4]{ThomasonTrobaugh}. Tensoring its exact sequence with the
perfect coefficient category, using \cite[Lemmas~3.12 and B.2]{LT}
and the preceding supported-kernel identification, gives the
coefficient version. Nonconnective $K$-theory sends this exact sequence
to a fiber sequence by \cite[Definition~8.1, Theorems~8.7 and 9.8,
and Section~9.3]{BGT}.

\subsection{Descent for finite surjections}

The functor $\cL$ is now a nilinvariant cdh sheaf computing actual
rational supported $K$-theory. To apply it to the finite map in
Theorem~\ref{ds:descent}, we need descent for finite surjections, which
need not be flat. Rational transfers give the finite flat case;
flattening by blowup and induction on closed subsets will give the
general case.

\begin{lemma}
\label{ds:finite-flat}
The functor $\cL$ satisfies ordinary \v{C}ech descent for finite
faithfully flat covers, universally under base change.
\end{lemma}

\begin{proof}
Let $f:X\to Y$ be finite faithfully flat. Since the schemes are
Noetherian, $f_*\mathcal O_X$ is a vector bundle of positive rank.
Restriction of scalars preserves perfect coefficient modules and
their supports: locally the scalar-extended coefficient algebra is
a finite projective module over the original one. It therefore
defines a transfer. Finite-flat base change and the projection
formula give natural transformations of presheaves on $Y$-schemes
\[
 \cL\xrightarrow{f^*}\cL_X
       \xrightarrow{\operatorname{tr}_f}\cL,
 \qquad
 \cL_X(W)=\cL(W\times_Y X),
\]
whose composite is tensoring with the pullback of
$f_*\mathcal O_X$.

Work first on a Zariski open of $Y$ on which that vector bundle is
free of rank $a>0$, and fix a trivialization there. Additivity of
$K$-theory identifies this composite with multiplication by $a$.
The same trivialization pulls back on every $Y$-scheme, so this
identity is natural on the whole presheaf diagram. Rationally,
$a^{-1}\operatorname{tr}_f$ makes $\cL$ a retract of $\cL_X$.

After base change to $X$, the augmented \v{C}ech nerve of $f$ has
a section and hence an extra degeneracy. Its image under any
presheaf has split descent. Thus $\cL_X$ satisfies descent for
$f$. The descent comparison for $\cL$ is a retract of this
equivalence, and is itself an equivalence. Finally, both its source
and its totalization are Zariski sheaves. The conclusion therefore
glues over the trivializing cover of $Y$. All constructions commute
with base change.
\end{proof}

\begin{lemma}
\label{ds:finite-surjection}
The functor $\cL$ satisfies ordinary \v{C}ech descent for every
finite surjection of finite-type $A$-schemes.
\end{lemma}

\begin{proof}
Let $\tau$ be the topology generated by cdh covers and finite
faithfully flat covers. These cover families are stable under base
change. By Lemma~\ref{ds:finite-flat} and cdh descent, $\cL$ is a
$\tau$-sheaf. We show that every finite surjection $X\to Y$ is a
$\tau$-cover.

Reductions and finite closed decompositions are cdh-local
operations. For clarity, if a reduced scheme is a union
$Y=Y_1\cup Y_2$ of reduced closed subschemes, the proper map
$Y_1\to Y$ is an isomorphism off $Y_2$ and gives the abstract
blowup square with intersection $Y_1\cap Y_2$. This reduces closed
unions to their pieces. We may consequently reduce to integral
$Y$, with induction on dimension for proper closed subsets.

Choose a reduced irreducible component $X_0$ of $X$ dominating
$Y$. The map $X_0\to Y$ is finite and surjective. On a dense open
$U\subset Y$ it is finite flat of positive rank. Flattening by
blowup, applied to $\mathcal O_{X_0}$, gives a blowup
$Y'\to Y$, an isomorphism over $U$, for which the strict transform
$X'_0\to Y'$ is flat and finitely presented
\cite[Tag~0815]{Stacks}. The strict transform is a closed subscheme
of $X_0\times_Y Y'$, so it is finite over $Y'$. The blowup $Y'$ is
integral, and the image of $X'_0$ contains its generic point.
Finiteness makes that image closed. Thus $X'_0\to Y'$ is
surjective, and is a finite faithfully flat cover factoring through
$X\times_Y Y'$.

Let $Z\subsetneq Y$ be the blowup center. The family
$\{Y'\to Y,Z\to Y\}$ is a cdh cover. By induction on dimension,
the finite surjection $X\times_Y Z\to Z$ is a $\tau$-cover.
Together with $X'_0\to Y'$, these maps give a $\tau$-covering
family of $Y$ that factors through $X$. Hence $X\to Y$ itself
generates a covering sieve. The induction starts in dimension zero,
where a reduced integral base is a field and any nonzero finite
algebra is finite faithfully flat. This proves the assertion, and
therefore the required \v{C}ech descent.
\end{proof}

The argument uses actual retracts and sheaf limits; it requires
neither hyperdescent nor convergence of an unbounded descent
spectral sequence.

\subsection{The geometric orbit argument}

We have established descent for finite surjections. It remains to use
geometric transitivity to identify the homotopy $G$-invariants with the
value on the base. After arranging a section and reducing the schemes,
we compare its translates away from their coincidence loci and use
induction along those loci.

\begin{proof}[Proof of Theorem~\ref{ds:descent}]
Write $X=\Spec T$ and $S=\Spec A$. On finite-type $S$-schemes
consider the natural comparison
\begin{equation}
 \cL(Y)\longrightarrow
 \cL(X\times_S Y)^{hG}.
 \label{ds:comparison}
\end{equation}
Both sides are sheaves for the topology used in
Lemma~\ref{ds:finite-surjection}: base change preserves the cover
families, and homotopy fixed points commute with limits. Applying
finite-surjection descent to $X\to S$, it suffices to prove
\eqref{ds:comparison} on the nonaugmented terms of its \v{C}ech
nerve. Over each such term, the pulled-back $X$-cover admits a
section, given by a projection to one of its $X$-coordinates. The
geometric transitivity hypothesis is stable under base change.
We are therefore reduced to covers admitting a section.

We prove this case by induction on the dimension of $Y$.
Nilinvariance and the closed-union squares described in the proof
of Lemma~\ref{ds:finite-surjection} reduce the base to an integral
reduced scheme. Indeed both sides of \eqref{ds:comparison} have
these descent properties; the component reduction uses finitely
many closed pieces and lower-dimensional intersections.

Let $s:Y\to X\times_S Y$ be a section. Its finitely many
$G$-translates cover the source on points. To see this, extend the
residue field at any point of $Y$ to an algebraic closure and use
transitivity on that geometric fiber. Each translate is a closed
immersion, because the source is separated over $Y$, and each
factors through the reduction
\[
 V=(X\times_S Y)_{\mathrm{red}}.
\]
Retain the distinct translated sections and index them by a
transitive finite $G$-set $I$. Two distinct sections have a proper
closed coincidence locus in $Y$: their equalizer is closed, and
agreement at the generic point of integral reduced $Y$ would imply
agreement everywhere. Let $Z\subsetneq Y$ be the reduced union
of these coincidence loci, with $Z=\varnothing$ if there is only
one section.

The map
\[
 q:W=\coprod_{i\in I}Y\longrightarrow V
\]
is finite and surjective. Over $Y\setminus Z$ the section images
are pairwise disjoint and cover the reduced scheme $V$. Their
disjoint union is therefore isomorphic to $V$ there. Consequently
$q$ is a proper modification and the square
\[
 \begin{array}{ccc}
 W\times_Y Z&\longrightarrow&W\\
 \big\downarrow&&\big\downarrow\\
 V\times_Y Z&\longrightarrow&V
 \end{array}
\]
is an abstract blowup square. It is $G$-equivariant.

The $G$-action on $W$ merely permutes its copies of $Y$.
The canonical restriction maps give a $G$-equivariant equivalence
$\cL(W)\simeq\prod_{i\in I}\cL(Y)$.
A stabilizer of one section acts identically on that copy, including
on the coefficient algebra pulled back from $Y$. This is a
canonical identity of coefficient-bearing $Y$-schemes, so the
action is the ordinary permutation action, with coherently trivial
stabilizer action. The same statement holds after restriction to $Z$.
Since $\cL$ is rational, averaging over a finite group is exact.
It follows that
\begin{equation}
 \cL(W)^{hG}\simeq\cL(Y),
 \qquad
 \cL(W\times_Y Z)^{hG}\simeq\cL(Z).
 \label{ds:permutation-invariants}
\end{equation}
For example, if $I=G/G_s$, the first fixed-point spectrum is
$\cL(Y)^{hG_s}$ for the trivial $G_s$-action, and is $\cL(Y)$.
This identity holds without boundedness: the trivial
$\Q[G_s]$-module $\Q$ is projective, so derived invariants are
ordinary invariants even on unbounded rational complexes.

Apply $\cL$ to the abstract blowup square and take homotopy
$G$-invariants. The inductive hypothesis over $Z$ identifies
\[
 \cL(V\times_Y Z)^{hG}\simeq\cL(Z).
\]
Here nilinvariance allows one to pass between $V\times_Y Z$ and
the corresponding reduced or unreduced pullback of the original
cover. Combining this with \eqref{ds:permutation-invariants}, the
pullback square gives
\[
 \cL(V)^{hG}
 \simeq\cL(Y)\mathop{\times}_{\cL(Z)}\cL(Z)
 \simeq\cL(Y).
\]
Nilinvariance restores $X\times_S Y$ in place of $V$. In
dimension zero there are no proper nonempty coincidence loci, and
the same proof is simply the calculation for a disjoint union of
section copies. This starts the induction and proves
\eqref{ds:comparison}.

Taking $Y=S$ and using \eqref{ds:support-identification} yields
\eqref{ds:spectral-descent}. Finally, exact rational invariants give
\[
 \pi_0\Bigl(\bigl(\K^J(H\otimes_A^{\mathbf L}T)_\Q\bigr)^{hG}\Bigr)
   =\bigl(K_0^J(H\otimes_A^{\mathbf L}T)\otimes\Q\bigr)^G,
\]
which proves \eqref{ds:kzero-descent}, in fact as an isomorphism.
\end{proof}

\subsection{Comparison of ordinary and normalized Euler lengths}

We now isolate the consequence needed before introducing asymptotic
categories. Let $(A,\m_A,k)$ be a complete regular local
$\Z_{(p)}$-algebra of dimension $e>0$, with residue characteristic
$p$, and let $z_1,\ldots,z_e$ be regular parameters. Let $T$ be a
finite local domain extension of $A$ with the same residue field.
Put $r=\rank_A T$. Suppose that $C$ is a $T$-algebra equipped with
an additive nonnegative length $\lambda$ on $C$-modules annihilated
by powers of $\m_A$, with $\lambda(0)=0$. Infinite length is
allowed for general modules. Assume that, for every $v\geq1$,
\begin{equation}
 \lambda H_i(z_1^v,\ldots,z_e^v;C)=0\quad(i>0),
 \qquad
 \lambda\bigl(C/(z_1^v,\ldots,z_e^v)C\bigr)=v^e r.
 \label{ds:length-contract}
\end{equation}
Theorem~\ref{len:input} supplies exactly these hypotheses in our
applications, since the multiplicity of the base parameters on
$T$ is its generic rank over the regular base.
The next lemma is the formal version of
Corollary~\ref{len:finite-length-base-change}, using only
\eqref{ds:length-contract}.

\begin{lemma}
\label{ds:finite-length-basechange}
For every finite-length $A$-module $V$,
\begin{equation}
 \lambda\Tor_i^A(V,C)=0\quad(i>0),
 \qquad
 \lambda(V\otimes_A C)=r\length_A(V).
 \label{ds:finite-length-formula}
\end{equation}
All these lengths are finite.
\end{lemma}

\begin{proof}
The Koszul resolution on $z_1,\ldots,z_e$ resolves $k$ over
$A$, so \eqref{ds:length-contract} with $v=1$ proves the
assertion for $V=k$. Induct on the length of $V$ using a
composition series. In the long exact Tor sequence associated to
a short exact sequence with quotient $k$, nonnegativity and
additivity show that every positive Tor term has length zero.
The image of the connecting map into degree zero also has length
zero. Additivity in degree zero then gives the asserted factor
$r$. This simultaneously proves finiteness at every step.
\end{proof}

If a finite group $G$ acts on $T$ over $A$, then for $g\in G$ scalar
twisting means $g^*P=P\otimes_{T,g}T$,
where the left $T$-action on the last factor is through $g$.
On a free complex this applies $g$ to every differential matrix entry.

\begin{lemma}[Ordinary norm comparison]
\label{ds:ordinary-norm}
Assume, in addition, that a finite group $G$ acts on $T$ over $A$
and is transitive on every geometric fiber of $\Spec T\to\Spec A$.
For a perfect $T$-complex $P$ supported on $V(\m_A T)$, the
quantities
\[
 \chi_T(P)=\sum_i(-1)^i\length_T H_i(P),
 \qquad
 \chi_C(P)=\sum_i(-1)^i
           \lambda H_i(P\otimes_T^{\mathbf L}C)
\]
are finite and define homomorphisms on $K_0^{\m_A}(T)$.
They agree on every norm sum:
\begin{equation}
 \sum_{g\in G}\chi_T(g^*P)
   =\sum_{g\in G}\chi_C(g^*P).
 \label{ds:ordinary-norm-identity}
\end{equation}
More generally, the two homomorphisms agree on the image of
$K_0^{\m_A}(A)$ and on every $G$-invariant rational class.
\end{lemma}

\begin{proof}
We first establish finiteness and additivity, then compare the
evaluations on pullbacks before applying descent to the norm class.
Represent $P$ by a bounded finite free $T$-complex. Its ordinary
homology is finite and supported at the closed point, hence has
finite length. For the $C$-evaluation, choose a common $v$ such
that each $z_j^v$ acts nullhomotopically on $P$. This is possible
by the support criterion above. Set
$K_v=K(z_1^v,\ldots,z_e^v;T)$. Successively splitting the cones
of these nullhomotopic actions gives
\begin{equation}
 K_v\otimes_T P\otimes_T C
 \simeq
 \bigoplus_{j=0}^e
       (P\otimes_T C)[j]^{\oplus\binom ej}.
 \label{ds:koszul-splitting}
\end{equation}
The homotopies may be tensored with $C$ and with the other Koszul
factors, so this splitting respects the module structure on which
$\lambda$ is defined.

In the opposite order, taking Koszul homology first gives a finite
spectral sequence whose terms are finite direct sums of
$H_j(z_1^v,\ldots,z_e^v;C)$. Their lengths are finite by
\eqref{ds:length-contract}. Hence every homology module on the
left side of \eqref{ds:koszul-splitting} has finite length.
The splitting and nonnegativity then give finiteness for every
$H_i(P\otimes_T C)$. The homology modules and spectral-sequence terms
to which $\lambda$ is applied are annihilated by fixed powers of
$\m_A$. The Euler sums are therefore defined, and the long exact
homology sequence of a triangle proves their additivity.

Let $Q$ now be a perfect $A$-complex supported at $\m_A$.
Its homology modules have finite length. The finite hyper-Tor
spectral sequence for $Q\otimes_A^{\mathbf L}C$, together with
Lemma~\ref{ds:finite-length-basechange}, gives
\begin{equation}
 \chi_C(Q\otimes_A^{\mathbf L}T)
       =r\sum_j(-1)^j\length_A H_j(Q).
 \label{ds:normalized-base-euler}
\end{equation}
Indeed every positive Tor row has length zero and the zeroth row
has the displayed scaled lengths.

For the ordinary evaluation, choose a bounded finite free
$A$-resolution $L_\bullet$ of $T$, possible because $A$ is
regular. Its alternating rank is $r$, as is seen after tensoring
with the fraction field of $A$. For every finite-length
$A$-module $V$, the finite complex $V\otimes_A L_\bullet$
therefore gives
\[
 \sum_i(-1)^i\length_A\Tor_i^A(V,T)
   =\sum_i(-1)^i\rank_A(L_i)\length_A(V)
   =r\length_A(V).
\]
The hyper-Tor spectral sequence for
$Q\otimes_A^{\mathbf L}T$ consequently yields
\begin{equation}
 \chi_T(Q\otimes_A^{\mathbf L}T)
       =r\sum_j(-1)^j\length_A H_j(Q).
 \label{ds:ordinary-base-euler}
\end{equation}
We have used that finite-length $T$-modules have the same
$A$-length and $T$-length, because both residue fields are $k$.
Equations~\eqref{ds:normalized-base-euler} and
\eqref{ds:ordinary-base-euler} prove agreement on classes pulled
back from $A$.

Apply Theorem~\ref{ds:descent} with $H=A$ and $J=\m_A$.
Every invariant class in $K_0^{\m_A}(T)\otimes\Q$ is the
pullback of a rational class downstairs. Extending the two Euler
homomorphisms $\Q$-linearly, they therefore agree on that
invariant subspace. The class $\sum_{g\in G}[g^*P]$ is
invariant, which gives \eqref{ds:ordinary-norm-identity}.
\end{proof}

\begin{remark}
\label{ds:no-length-invariance}
No invariance of $\lambda$ under $G$ is asserted or needed.
The equality is for the sum of the conjugate classes, not for each
individual class. This distinction allows the two length algebras
used later to be chosen independently.
\end{remark}
